\documentclass[11pt]{article}
\usepackage[margin=1in]{geometry}
\usepackage[hidelinks]{hyperref}
\usepackage{tabularx}
\usepackage{booktabs}
\usepackage{enumitem}
\usepackage{placeins}
\usepackage{url}

\title{Direct and Anchor-Gated Polyp Sizing from Monocular Colonoscopy Video\\
\large Literature Review and Research Direction}
\author{Smart Wattanapornmongkol\thanks{\raggedright Handoff contact: student ID 6738249221; phone 0639390653; \href{https://line.me/ti/p/nqlvyY38hD}{LINE}.\\Project repository: \href{https://github.com/smartwhatt/Colon-AF-SfMLearner}{github.com/smartwhatt/Colon-AF-SfMLearner}.}}
\date{23 July 2026}

\begin{document}
\maketitle

\begin{abstract}
Polyp size is clinically consequential, but estimating it from a monocular
colonoscopy video is not the same task as estimating relative depth. This
review introduces the clinical and camera context needed by machine-learning
researchers, surveys direct and reference-based size-estimation approaches,
and identifies their shared limitations. Wide-angle optics, near-light
illumination, non-rigid tissue, quality variation, and missing metric scale
make single-frame pixel measurements unreliable. Existing approaches use
learned RGB/depth features, temporal aggregation, calibrated optical priors,
or visible references such as a waterjet or snare. The recommended research
program is an evaluation-first, short-video direct-sizing model: track a
lesion, filter for usable views, predict an ordinal or continuous endpoint,
aggregate at polyp level, and emit millimetres only when a validated scale
anchor is available. In all other cases, the model should return a calibrated
ordinal estimate or abstain rather than imply unsupported metric precision.
\end{abstract}

\section{Scope and terminology}

This is a methods review for research planning, not clinical advice. A
\emph{polyp} is a lesion on the inner lining of the large bowel. During
colonoscopy, an endoscopist inspects the lesion, may remove it, and records an
estimated diameter. In research, $\leq 5$~mm, 6--9~mm, and $\geq 10$~mm are
useful ordinal ranges; they should not be treated as a substitute for clinical
guidelines or an automatic treatment policy.

The terms \emph{metric size}, \emph{relative depth}, and \emph{ordinal size}
must remain separate. Metric size is a physical dimension in millimetres.
Relative depth is a geometrically consistent shape up to a global scale.
Ordinal size predicts a clinically useful range without asserting millimetres.
This review distinguishes methods that output a number from methods that have
validated a transferable scale channel.

\section{Clinical and camera context}

\subsection{What is measured in practice}

Labels called \texttt{size\_mm} are not interchangeable. A fresh-specimen
caliper records an excised object that may have changed through handling;
pathology records can reflect fixation or fragmentation; an endoscopist's
visual estimate is clinically useful but not independent metric ground truth;
and a snare or forceps ratio depends on tool opening, visibility, pose, and
whether the tool lies at the lesion depth. Phantom geometry offers controlled
ground truth but does not reproduce clinical tissue or workflow. Every
evaluation manifest should record label provenance, measurement protocol,
scope model, site, operator, and whether an in-scene reference is present.

\subsection{Why a colonoscopy image is difficult}

A colonoscope is a wide-angle endoscope with illumination close to the lens.
Its imagery differs from a fixed, well-calibrated pinhole camera in several
ways: field-of-view distortion and image processing vary by scope; brightness
changes with distance, orientation, exposure, white balance, and tissue
reflectance; tissue deforms under insufflation, contact, fluid, peristalsis,
and scope pressure; and procedure video includes blur, bubbles, debris,
instruments, treatment, and occlusion. A lesion clip is therefore a selected
subset of an uneven, quality-variable procedure video, not simply many
independent still images.

\subsection{Metric-scale ambiguity}

For a perspective camera, $x=fX/Z$. Multiplying $X$, $Y$, $Z$, and camera
translation by the same positive constant leaves image projection unchanged.
Monocular structure-from-motion and photometric self-supervision can recover
relative geometry without recovering its physical scale. Median scaling can
evaluate relative shape, but never proves millimetres. Metric output needs an
independent scale cue---for example a known visible reference, a calibrated
lens-and-procedure prior that is validated for the target setting, or a
scope-specific empirical calibration.

\section{Methods in the literature}

Table~\ref{tab:methods} groups representative approaches by the evidence that
links visual features to size. The table is deliberately not a leaderboard:
cohorts, label protocols, endpoint definitions, and split units differ too
substantially for the reported numbers to be ranked across papers.

\begin{table}[t]
\centering
\footnotesize
\setlength{\tabcolsep}{5pt}
\begin{tabularx}{\textwidth}{@{}p{0.28\textwidth}X@{}}
\toprule
\textbf{Approach and representative work} & \textbf{Endpoint, scale evidence, and planning implication} \\
\midrule
RGB plus self-supervised depth and temporal pooling: Itoh et al. (BseNet) \cite{itoh2018bsenet} & Binary $<10$ versus $\geq10$~mm classification; depth is an auxiliary learned cue rather than an explicit physical reference. Direct clinical endpoint and temporal aggregation are feasible, but binary accuracy does not validate metric scale or cross-scope transfer. \\

Attention-gated multi-cue classifiers: Roffo et al. \cite{roffo2024hag} (preprint) & RGB, depth, and location features for binary and three-class sizing. Combining cues is plausible; evaluate at lesion level and enforce site/scope/entity-disjoint splits before drawing general conclusions. \\

Reference-free direct regression: Yii et al. \cite{yii2025referencefree} & Segmentation and depth cues with direct millimetre regression, evaluated under a single-scope, snare-anchored protocol. A direct regression head is viable, but the protocol's anchor and hardware context are part of the result. \\

Geometry-aware direct regression: Lawrence et al. (MPSE) \cite{lawrence2026mpse} & Per-frame RGB/mask/pseudodepth regression with polyp-level aggregation. Aggregation should be designed around a tracked lesion rather than chosen post hoc; small or imbalanced cohorts require uncertainty reporting. \\

Calibrated optical prior: ENDOANGEL \cite{wang2024endoangel} & A real-time deep-learning system that incorporates a lens-parameter calibration channel. Treat calibration as a target-environment requirement and validate it prospectively across hardware and sites. \\

Visible physical reference: Poseidon \cite{sudarevic2023poseidon} & The auxiliary waterjet serves as a reference object for size measurement. A deliberately visible anchor is often the clearest path from pixels to millimetres, but introduces workflow and visibility failure modes. \\
\bottomrule
\end{tabularx}
\caption{Representative methods and the scale evidence they use. Preprint
status is stated explicitly; its findings should not be treated as mature
clinical evidence.}
\label{tab:methods}
\end{table}

Early direct approaches demonstrate why a reconstruction-only framing is too
narrow. BseNet predicts self-supervised relative depth, combines it with RGB,
and pools temporal information for a binary size decision \cite{itoh2018bsenet}.
This makes the output closer to a clinical question, but does not resolve
metric identifiability. More recent direct models use segmentation, estimated
depth, and learned fusion to regress a diameter or class
\cite{yii2025referencefree,lawrence2026mpse}. Their central contribution is
the endpoint formulation; their measurement claims still depend on the label
and scale protocol.

Reference-based work makes that protocol explicit. ENDOANGEL uses calibrated
lens parameters in a multicentre prospective validation \cite{wang2024endoangel}.
Poseidon uses the waterjet as a physical in-frame reference
\cite{sudarevic2023poseidon}. These strategies carry different operational
costs, but both expose a measurement channel that a generic unanchored video
cannot guarantee.

Temporal representation learning is also relevant. Parolari et al. study
contrastive learning under noisy temporal self-supervision and evaluate
downstream colonoscopy-video tasks including size prediction
\cite{parolari2026temporal}; this is a 2026 preprint rather than clinical
validation. Its useful lesson for a sizing project is to model a lesion track,
not merely to average arbitrary adjacent frames.

\FloatBarrier
\section{Failure modes that should shape evaluation}

The most common risk is false precision: a model returns millimetres because
its head is trained to regress a number, while its evaluation does not prove
that scale transfers across scopes, sites, or label sources. Other predictable
failure modes are frame-level leakage from the same lesion into train and test;
mixed label provenance; tool-anchor occlusion or orientation error; blurred or
near-contact views; under-representation of large polyps; and cherry-picking
the best frame after observing labels. A dataset split should be at least
polyp- and procedure-disjoint, and when data permit, site-, scope-, and
operator-aware. Report performance by clinically relevant range, label source,
anchor visibility, scope model, and image-quality stratum.

For a continuous metric endpoint, report MAE together with agreement analysis,
calibration, confidence intervals at the lesion level, and an abstention curve.
For an ordinal endpoint, report confusion across clinically defined ranges,
calibration of class probabilities, and selective risk as low-confidence cases
are deferred. A fixed qualitative panel should include specular highlights,
blur, bubbles, instruments, partial lesions, low insufflation, and difficult
viewing angles. Do not use median-scaled depth as evidence of metric accuracy.

\section{Recommended research program}

The recommended path is an \emph{anchor-gated short-video direct-sizing}
track, run beside the project’s existing geometry guardrails rather than
assuming the depth-pose stack is the sole solution.

\begin{enumerate}[leftmargin=*]
  \item \textbf{Create an evaluation manifest first.} Deduplicate at lesion
  and procedure level; preserve label provenance; mark scope/site/operator;
  record anchor type and visibility; and reserve leakage-checked held-out
  lesions before selecting an architecture.
  \item \textbf{Form lesion clips.} Detect or segment the lesion, track it
  through a short clip, and reject frames with blur, occlusion, poor lesion
  coverage, excessive specularity, or unreliable tool/anchor visibility.
  \item \textbf{Use two linked outputs.} Train an ordinal head for all
  quality-approved clips. Train a continuous size head only for examples where
  the scale evidence is valid; predict uncertainty for both outputs.
  \item \textbf{Aggregate at polyp level.} Combine frame predictions using a
  pre-specified robust rule, attention, or learned temporal model. The unit of
  selection and reporting should be the lesion, not the most favorable frame.
  \item \textbf{Gate metric output.} Emit millimetres only when anchor
  detection, calibration, image quality, and uncertainty all pass. Otherwise
  output the ordinal distribution or abstain.
  \item \textbf{Compare against useful baselines.} Include a single-frame RGB
  ordinal model, a clip model without depth features, an optional depth-feature
  model, and a reference-based baseline when anchors are available.
\end{enumerate}

This approach is deliberately conservative. It permits a model to be useful in
the many views where millimetres are not defensible, while making the conditions
for a metric claim inspectable. The project’s self-supervised depth model can
enter as a feature or a quality signal, but it should be evaluated as an
incremental contribution rather than assumed to solve scale.

\section{Conclusion}

Direct polyp sizing from video is a legitimate endpoint, but direct regression
alone does not solve monocular scale ambiguity. The strongest next experiment
is a leakage-controlled, polyp-level benchmark with explicit anchor and label
provenance, followed by a short-video model that separates ordinal inference
from anchor-supported metric output. Such a design is more clinically legible,
more robust to missing scale evidence, and easier to falsify than a claim that
unanchored monocular reconstruction recovers millimetres.

\begingroup
\scriptsize
\bibliographystyle{plain}
\bibliography{references}
\endgroup
\end{document}
